% Options for packages loaded elsewhere
% Options for packages loaded elsewhere
\PassOptionsToPackage{unicode}{hyperref}
\PassOptionsToPackage{hyphens}{url}
\PassOptionsToPackage{dvipsnames,svgnames,x11names}{xcolor}
%
\documentclass[
  letterpaper,
  DIV=11,
  numbers=noendperiod]{scrreprt}
\usepackage{xcolor}
\usepackage{amsmath,amssymb}
\setcounter{secnumdepth}{5}
\usepackage{iftex}
\ifPDFTeX
  \usepackage[T1]{fontenc}
  \usepackage[utf8]{inputenc}
  \usepackage{textcomp} % provide euro and other symbols
\else % if luatex or xetex
  \usepackage{unicode-math} % this also loads fontspec
  \defaultfontfeatures{Scale=MatchLowercase}
  \defaultfontfeatures[\rmfamily]{Ligatures=TeX,Scale=1}
\fi
\usepackage{lmodern}
\ifPDFTeX\else
  % xetex/luatex font selection
\fi
% Use upquote if available, for straight quotes in verbatim environments
\IfFileExists{upquote.sty}{\usepackage{upquote}}{}
\IfFileExists{microtype.sty}{% use microtype if available
  \usepackage[]{microtype}
  \UseMicrotypeSet[protrusion]{basicmath} % disable protrusion for tt fonts
}{}
\makeatletter
\@ifundefined{KOMAClassName}{% if non-KOMA class
  \IfFileExists{parskip.sty}{%
    \usepackage{parskip}
  }{% else
    \setlength{\parindent}{0pt}
    \setlength{\parskip}{6pt plus 2pt minus 1pt}}
}{% if KOMA class
  \KOMAoptions{parskip=half}}
\makeatother
% Make \paragraph and \subparagraph free-standing
\makeatletter
\ifx\paragraph\undefined\else
  \let\oldparagraph\paragraph
  \renewcommand{\paragraph}{
    \@ifstar
      \xxxParagraphStar
      \xxxParagraphNoStar
  }
  \newcommand{\xxxParagraphStar}[1]{\oldparagraph*{#1}\mbox{}}
  \newcommand{\xxxParagraphNoStar}[1]{\oldparagraph{#1}\mbox{}}
\fi
\ifx\subparagraph\undefined\else
  \let\oldsubparagraph\subparagraph
  \renewcommand{\subparagraph}{
    \@ifstar
      \xxxSubParagraphStar
      \xxxSubParagraphNoStar
  }
  \newcommand{\xxxSubParagraphStar}[1]{\oldsubparagraph*{#1}\mbox{}}
  \newcommand{\xxxSubParagraphNoStar}[1]{\oldsubparagraph{#1}\mbox{}}
\fi
\makeatother


\usepackage{longtable,booktabs,array}
\usepackage{calc} % for calculating minipage widths
% Correct order of tables after \paragraph or \subparagraph
\usepackage{etoolbox}
\makeatletter
\patchcmd\longtable{\par}{\if@noskipsec\mbox{}\fi\par}{}{}
\makeatother
% Allow footnotes in longtable head/foot
\IfFileExists{footnotehyper.sty}{\usepackage{footnotehyper}}{\usepackage{footnote}}
\makesavenoteenv{longtable}
\usepackage{graphicx}
\makeatletter
\newsavebox\pandoc@box
\newcommand*\pandocbounded[1]{% scales image to fit in text height/width
  \sbox\pandoc@box{#1}%
  \Gscale@div\@tempa{\textheight}{\dimexpr\ht\pandoc@box+\dp\pandoc@box\relax}%
  \Gscale@div\@tempb{\linewidth}{\wd\pandoc@box}%
  \ifdim\@tempb\p@<\@tempa\p@\let\@tempa\@tempb\fi% select the smaller of both
  \ifdim\@tempa\p@<\p@\scalebox{\@tempa}{\usebox\pandoc@box}%
  \else\usebox{\pandoc@box}%
  \fi%
}
% Set default figure placement to htbp
\def\fps@figure{htbp}
\makeatother





\setlength{\emergencystretch}{3em} % prevent overfull lines

\providecommand{\tightlist}{%
  \setlength{\itemsep}{0pt}\setlength{\parskip}{0pt}}



 


\KOMAoption{captions}{tableheading}
\makeatletter
\@ifpackageloaded{bookmark}{}{\usepackage{bookmark}}
\makeatother
\makeatletter
\@ifpackageloaded{caption}{}{\usepackage{caption}}
\AtBeginDocument{%
\ifdefined\contentsname
  \renewcommand*\contentsname{Table of contents}
\else
  \newcommand\contentsname{Table of contents}
\fi
\ifdefined\listfigurename
  \renewcommand*\listfigurename{List of Figures}
\else
  \newcommand\listfigurename{List of Figures}
\fi
\ifdefined\listtablename
  \renewcommand*\listtablename{List of Tables}
\else
  \newcommand\listtablename{List of Tables}
\fi
\ifdefined\figurename
  \renewcommand*\figurename{Figure}
\else
  \newcommand\figurename{Figure}
\fi
\ifdefined\tablename
  \renewcommand*\tablename{Table}
\else
  \newcommand\tablename{Table}
\fi
}
\@ifpackageloaded{float}{}{\usepackage{float}}
\floatstyle{ruled}
\@ifundefined{c@chapter}{\newfloat{codelisting}{h}{lop}}{\newfloat{codelisting}{h}{lop}[chapter]}
\floatname{codelisting}{Listing}
\newcommand*\listoflistings{\listof{codelisting}{List of Listings}}
\makeatother
\makeatletter
\makeatother
\makeatletter
\@ifpackageloaded{caption}{}{\usepackage{caption}}
\@ifpackageloaded{subcaption}{}{\usepackage{subcaption}}
\makeatother
\usepackage{bookmark}
\IfFileExists{xurl.sty}{\usepackage{xurl}}{} % add URL line breaks if available
\urlstyle{same}
\hypersetup{
  pdftitle={Notebook: GenAI},
  pdfauthor={Angelika Merkel},
  colorlinks=true,
  linkcolor={blue},
  filecolor={Maroon},
  citecolor={Blue},
  urlcolor={Blue},
  pdfcreator={LaTeX via pandoc}}


\title{Notebook: GenAI}
\author{Angelika Merkel}
\date{2025-12-04}
\begin{document}
\maketitle

\renewcommand*\contentsname{Table of contents}
{
\hypersetup{linkcolor=}
\setcounter{tocdepth}{2}
\tableofcontents
}

\bookmarksetup{startatroot}

\chapter*{Preface}\label{preface}
\addcontentsline{toc}{chapter}{Preface}

\markboth{Preface}{Preface}

This is a notebook about my personal journey into generative AI, which I
started during a break from work in the fall and winter of 2025. Here's
the rough path I followed:

\begin{itemize}
\item
  \textbf{A short history of generative AI} -- I started by using
  \href{https://chatgpt.com}{ChatGPT} and
  \href{https://www.perplexity.ai}{Perplexity} to get some background:
  how we got here, which milestones really mattered, and who the main
  players and models are today.
\item
  \textbf{ChatGPT, the basics} -- I subscribed to ChatGPT Pro (22
  EUR/month) and worked through the beginner tutorials at the
  \href{https://academy.openai.com}{OpenAI Academy} (free). This was
  quick but surprisingly solid to build a foundation and understand what
  the tool can do (and what not) and how to use the interface.
\item
  \textbf{Prompt engineering} -- To learn what I call \emph{The art of
  talking to the machine}, I subscribed to
  \href{https://www.coursera.org}{Coursera} (42 EUR/month) and took some
  basic and advanced courses on prompt engineering by Jules White,
  including how to create custom AI assistants (Custom GPTs) and build
  trustworthy genAI. This is where I started to get some real insights
\item
  \textbf{Data analysis with ChatGPT} -- Also by Jules White, these
  courses go deeper into practical workflows, showing how to use ChatGPT
  (via Code Interpreter) together with external tools for real data
  analysis tasks (and real magic).
\item
  \textbf{Risks and Ethics of generative AI} - A complication of
  insights and resources on how to provide appropriate rail guards for
  generative AI, AI agents and agentic AI and how (in my opinion)
  generative AI should be used ethically
\end{itemize}

\bookmarksetup{startatroot}

\chapter{Prompt Engineering}\label{prompt-engineering}

Based on my notes from the coursera courses '\emph{Prompt Engineering
with ChatGTP'}, \emph{`Advanced Prompt Engineering'}, and \emph{`Build
Your CustomGPT'} by Jules White(Vanderbuilt University). Adapted with
ChatGTP.

\section{How LLMs Work (Conceptual + Technical
Basics)}\label{how-llms-work-conceptual-technical-basics}

\begin{itemize}
\item
  \textbf{Large Language Models (LLMs)} generate text by predicting the
  \emph{next token} based on patterns learned from training data.

  → They operate \textbf{autoregressively}: each generated token is
  appended to the context and used to predict the next one.
\item
  \textbf{ChatGPT} = \emph{Chat Generative Pre-trained Transformer}:

  \begin{itemize}
  \tightlist
  \item
    \textbf{Self-attention} computes weighted relationships across all
    tokens, allowing the model to focus on the most relevant parts of
    the context.
  \item
    \textbf{Positional encodings} inject information about token order,
    enabling the model to represent sequences rather than bags of words.
  \end{itemize}
\item
  \textbf{Internal processing pipeline}:

  \begin{itemize}
  \tightlist
  \item
    Text → \textbf{tokenization} → \textbf{numerical embeddings} →
    stacked \textbf{transformer layers} → probability distribution over
    the next token.
  \end{itemize}
\item
  \textbf{Probability, not understanding}:

  \begin{itemize}
  \tightlist
  \item
    The model does not ``know'' facts; it estimates which continuation
    is most likely given the context and its training distribution.
  \end{itemize}
\item
  \textbf{Training objective}:

  \begin{itemize}
  \tightlist
  \item
    Minimize \textbf{cross-entropy loss}, pushing predicted token
    probabilities closer to the tokens observed in large-scale text
    corpora.
  \end{itemize}
\item
  \textbf{Determinism vs.~randomness}:

  \begin{itemize}
  \tightlist
  \item
    Sampling parameters (e.g.~temperature, top-k/top-p) influence how
    deterministic or creative the output appears.
  \end{itemize}
\end{itemize}

\textbf{Example:}\\
If you write \emph{``The mitochondria is the \ldots{}''}, the model
predicts \emph{``powerhouse of the cell''} because this phrase has a
high probability in the training data.

\begin{center}\rule{0.5\linewidth}{0.5pt}\end{center}

\section{Prompt Engineering (Purpose and
Practice)}\label{prompt-engineering-purpose-and-practice}

\begin{quote}
\emph{``The art of talking to the machine.''}
\end{quote}

\textbf{Prompt engineering aims to:}

\begin{itemize}
\tightlist
\item
  Produce \textbf{reliable}, \textbf{repeatable}, and \textbf{useful}
  outputs despite inherent nondeterminism.
\item
  Increase \textbf{predictability}, \textbf{consistency}, and
  \textbf{task alignment}.
\item
  Clearly encode \textbf{task intent}, \textbf{context},
  \textbf{constraints}, and \textbf{output structure}.
\item
  Reduce ambiguity and lower the risk of hallucinations.
\end{itemize}

\textbf{Example 1 (general):}

Instead of:\\
\textgreater{} ``Explain CRISPR.''

Better:\\
\textgreater{} ``Act as a molecular biologist. Explain CRISPR to a
motivated high-school student in 3 bullet points.''

\textbf{Example 2 (technical):}

Instead of:\\
\textgreater{} ``Explain PCA.''

Better:\\
\textgreater{} ``Provide a mathematically grounded explanation of PCA
for an R programmer. Include eigenvectors, variance decomposition, and a
3-line R code example.''

\bookmarksetup{startatroot}

\chapter{Prompt Pattern (practical)}\label{prompt-pattern-practical}

-- A practical, example-driven overview.

\hfill\break

\section{Core Prompt Patterns}\label{core-prompt-patterns}

\subsection{Persona Pattern}\label{persona-pattern}

Helps control tone, expertise, cultural assumptions.

\textbf{Example:}

``Act as a senior bioinformatician experienced in single-cell
epigenomics. Provide actionable suggestions.''

\subsection{Audience Persona Pattern}\label{audience-persona-pattern}

Tell GPT \emph{who you are} so explanations match your background.

\textbf{Example:}

``Assume I am new to machine learning but familiar with R and
statistics. Explain what a transformer is.''

\subsection{Question Refinement}\label{question-refinement}

GPT improves your question. Useful for technical tasks where question
quality matters.

\textbf{Example:}

``When I ask you something unclear, suggest a sharper version first and
ask if I want to use it.''

\subsection{Few-shot Prompting}\label{few-shot-prompting}

You show examples (introduce training data) → GPT follows your pattern.
Useful for labeling, transformations, extraction tasks.

\textbf{Example:}

Input: ATGCGT

Output: reversed sequence: TGCGTA

\hfill\break
Prompt:

Input: AACCAT

Output:

GPT continues with the intended pattern.

\subsection{Flipped Interaction}\label{flipped-interaction}

GPT asks \emph{you} the questions.

\textbf{Example:}

``Ask me clarifying questions until you can design a workshop outline.
When you have enough information, present the final outline.''

\subsection{Chain of Thought}\label{chain-of-thought}

Explicit reasoning steps.

\textbf{Example:}

``Show your reasoning in short bullet points and end with a boxed final
answer.''

Helps when solving problems or planning.

\subsection{ReAct Prompting}\label{react-prompting}

Reason → take external action → reason → answer.

\textbf{Example:}

``Search PubMed for recent papers on DNA methylation in HSCs, summarize
top 3 findings, and explain how they relate to AML biology.''

\emph{(Assuming GPT has tools activated.)}

\subsection{Alternative Approaches}\label{alternative-approaches}

Ask for multiple solutions. (or strategy for approaches)

\textbf{Example:}

``Propose 5 different ways to create a metadata-capture workflow for a
small lab, ranked by feasibility.''

Great for brainstorming and problem solving.

\subsection{Evaluation Pattern}\label{evaluation-pattern}

Use examples or rubrics to grade model output.

\textbf{Example:}

``Here are 3 example abstracts. Score the following abstract for clarity
(1--5), structure (1--5), scientific accuracy (1--5). Provide
justification.''

\subsection{Game-play Pattern}\label{game-play-pattern}

Turn tasks into structured exercises. Define rules.

\textbf{Example:}

``We are going to play a game:\\
I describe a lab workflow, you generate 5 mistakes a junior researcher
might make, and how to fix them.''

\subsection{Recipe Pattern}\label{recipe-pattern}

Stepwise instructions with gaps.

\textbf{Example:}

``Fill in missing steps for this workflow:

\begin{enumerate}
\def\labelenumi{\arabic{enumi}.}
\item
\item
  Load FASTQ files
\item
\item
  Align using BWA
\item
\item
  Generate BAM summary''
\item
\end{enumerate}

\subsection{Template Pattern}\label{template-pattern}

Use placeholders → consistent output.

\textbf{Example:}

Mutation Summary

Gene: \textless GENE\textgreater{}

Variant: \textless VARIANT\textgreater{}

Impact: \textless IMPACT\textgreater{}

Evidence: \textless QUOTES FROM ARTICLE\textgreater{}

\section{Text Manipulation Patterns}\label{text-manipulation-patterns}

\subsection{Semantic Filter}\label{semantic-filter}

Remove or replace specific elements.

\textbf{Example:}

``Remove all dates from the text but change as little as possible.''

\subsection{Meta-language Pattern}\label{meta-language-pattern}

Define a micro-language the LLM should learn.

\textbf{Example:}

``When I write `QC-S1', interpret it as: run initial QC on sequencing
files and summarize FastQC results.''

\section{Verifier Pattern}\label{verifier-pattern}

\subsection{Cognitive Verifier
Pattern}\label{cognitive-verifier-pattern}

Meta-reasoning pattern where the model is instructed to generate an
answer, independently verify, critique, or test that answer and revise
if needed. Focuses on error detection, not explanation.

Typical prompt form:

``First provide an answer. Then verify it by checking assumptions, edge
cases, and alternative interpretations.'' ``Critically review your own
solution and correct any mistakes.''

What it does well - Reduces hallucinations - Improves factual and
logical reliability - Surfaces hidden assumptions - Excellent for
high-stakes outputs

What it is not - Not primarily about explaining reasoning to the user -
Verification may be concise or structured rather than narrative

\textbf{Example:}

Answer: X

Verification: - Check assumptions → valid - Alternative method → same
result - Edge cases → none violated

Conclusion: Answer stands.

\subsection{Fact Checking Pattern}\label{fact-checking-pattern}

Force verification.

Steps:

\begin{enumerate}
\def\labelenumi{\arabic{enumi}.}
\setcounter{enumi}{1}
\tightlist
\item
  Extract facts
\item
  Compare with source
\item
  Flag inconsistencies
\item
  Insert verified fact list at the end
\end{enumerate}

\textbf{Example:}

``Extract all factual statements from the text and list any that need
verification.''

\section{Workflows}\label{workflows}

\subsection{Working on smaller pieces}\label{working-on-smaller-pieces}

Outline Expansion Start small, expand in pieces.

\textbf{Example:}

``Here is the outline for a workshop. Expand only section 2. Keep the
rest unchanged.''

\subsection{Tail Generation}\label{tail-generation}

Add automatic ``post-processing.''

\textbf{Example:}

``When done summarizing, also produce a 3-item action list.''

\bookmarksetup{startatroot}

\chapter{Prompt Pattern (technical)}\label{prompt-pattern-technical}

Technical Summary (with Examples)

\hfill\break

\emph{A compact, engineering-oriented overview}

\section{Core Prompt Pattern}\label{core-prompt-pattern}

\subsection{Persona Pattern (Functional
Control)}\label{persona-pattern-functional-control}

Define computational role + constraints.

\textbf{Example:}

``Act as a senior data engineer. Output only SQL, no explanations.''

\subsection{Audience Persona Pattern (Adjust Explanation
Level)}\label{audience-persona-pattern-adjust-explanation-level}

Tells the model who the target reader is.

\textbf{Example:}

``Explain transformer attention to someone familiar with matrix
multiplication but new to NLP.''

\subsection{Flipped Interaction Pattern (Information
Acquisition)}\label{flipped-interaction-pattern-information-acquisition}

Useful when task requires parameters or architecture definition.

\textbf{Example:}

``Ask me all required questions to define a Snakemake workflow. Then
generate the final workflow.''

\subsection{Question Refinement Pattern (Error
Reduction)}\label{question-refinement-pattern-error-reduction}

GPT rewrites your question to reduce ambiguity.

\textbf{Example:}

``Before answering, propose a more precise version of my question
optimized for computational reproducibility.''

\subsection{Cognitive Verifier Pattern (Structured
Reasoning)}\label{cognitive-verifier-pattern-structured-reasoning}

Model decomposes problem → produces final synthesized answer.

Useful to check the model's reasoning. Improves the output.

\textbf{Example:}

``Break the question into technical sub-questions (≤5), answer them,
then produce the final integrated explanation.''

\subsection{Few-Shot Pattern (Pattern
Imitation)}\label{few-shot-pattern-pattern-imitation}

Provide input/output examples → model induces rule.

\textbf{Example:}

Input: ``ATGGCT''

Output: reverse\_complement: ``AGCCAT''

\hfill\break
Input: ``CGTAAT''

Output:

GPT generalizes rule across new inputs.

\subsection{Chain-of-Thought (CoT) Pattern (Explain
Reasoning)}\label{chain-of-thought-cot-pattern-explain-reasoning}

Model outputs intermediate reasoning (optional).

\textbf{Example:}

``Show your reasoning in short bullet points and end with a boxed final
answer.''

\subsection{ReAct Pattern (Reason + Tool
Use)}\label{react-pattern-reason-tool-use}

LLM reasons → takes actions (search, code execution) → final answer.

\textbf{Example:}

``Search the latest ENCODE metadata for a DNase-seq dataset, summarize
QC metrics, then evaluate quality.''

\subsection{Evaluation Pattern (Auto-Grading
Models)}\label{evaluation-pattern-auto-grading-models}

Provide criteria or examples → GPT scores candidates.

\textbf{Example:}

``Evaluate the following three pipeline descriptions for clarity (1--5),
modularity (1--5), reproducibility (1--5). Justify each score.''

\subsection{Template Pattern (Deterministic
Structure)}\label{template-pattern-deterministic-structure}

Force GPT to output in a machine-readable or human-readable schema.

\textbf{Example:}

\{

``task'': ``\textless TASK\textgreater{}'',

``inputs'': \["\<INPUT_1\>", "\<INPUT_2\>"\],

``steps'': \["\<STEP_1\>", "\<STEP_2\>"\],

``expected\_output'': ``\textless OUTPUT\textgreater{}''

\}

\subsection{Meta-Language Pattern (Define Custom
Syntax)}\label{meta-language-pattern-define-custom-syntax}

You define a symbolic shorthand → model learns mapping.

\textbf{Example:}

``Interpret `QC-S1' as: run FastQC → summarize per-base content → output
3 metrics.''

\subsection{Recipe Pattern (Pipeline
Definition)}\label{recipe-pattern-pipeline-definition}

Fill missing steps in a workflow.

\textbf{Example:}

``Insert the two missing steps in this pipeline:

\begin{enumerate}
\def\labelenumi{\arabic{enumi}.}
\item
  Load FASTQ
\item
  Align with BWA
\item
  Generate BAM summary''
\end{enumerate}

\subsection{Alternative Approaches Pattern (Algorithmic
Exploration)}\label{alternative-approaches-pattern-algorithmic-exploration}

Generate multiple candidate solutions.

\textbf{Example:}

``Provide 4 alternative strategies for building a FAIR-compliant
metadata capture workflow. Rank by implementation cost.''

\section{Text \& Data Manipulation
Patterns}\label{text-data-manipulation-patterns}

\subsection{Semantic Filter}\label{semantic-filter-1}

Remove or transform only the targeted elements.

\textbf{Example:}

``Remove all patient identifiers but keep clinical meaning intact.''

\subsection{Tail Pattern
(Post-Processing)}\label{tail-pattern-post-processing}

Append additional required content automatically.

\textbf{Example:}

``After summarizing the paper, append a list of unresolved questions.''

\subsection{Fact-Checking Pattern}\label{fact-checking-pattern-1}

Extract facts → verify → append fact list.

\textbf{Example:}

``Extract all factual claims from the text and mark those that require
verification due to missing sources.''

\bookmarksetup{startatroot}

\chapter{CustomGPTs}\label{customgpts}

\section{Intro}\label{intro}

User-configured versions of ChatGPT that behave like a specialized
assistant for a specific task, domain, or workflow---without writing
code. Persistent behaviour.

Can be created through the ChatGTP app.

\begin{longtable}[]{@{}
  >{\raggedright\arraybackslash}p{(\linewidth - 4\tabcolsep) * \real{0.2500}}
  >{\raggedright\arraybackslash}p{(\linewidth - 4\tabcolsep) * \real{0.3194}}
  >{\raggedright\arraybackslash}p{(\linewidth - 4\tabcolsep) * \real{0.4306}}@{}}
\toprule\noalign{}
\begin{minipage}[b]{\linewidth}\raggedright
Aspect
\end{minipage} & \begin{minipage}[b]{\linewidth}\raggedright
ChatGPT
\end{minipage} & \begin{minipage}[b]{\linewidth}\raggedright
Custom GPT
\end{minipage} \\
\midrule\noalign{}
\endhead
\bottomrule\noalign{}
\endlastfoot
Purpose & General-purpose assistant & Task- or domain-specific
assistant \\
Instructions & Provided ad hoc per chat & Persistent, predefined
instructions \\
Knowledge & Generic model knowledge & Optional uploaded documents \&
references \\
Consistency & Varies between chats & Stable, repeatable behavior \\
Prompting effort & Repeated each session & One-time setup \\
Tool access & User-controlled per chat & Preconfigured tools \\
Hallucination risk & Higher (open-ended) & Lower (scoped +
constrained) \\
Best use case & Exploration, quick questions & Workflows, training,
institutional use \\
\end{longtable}

/

\textbf{Components:}

\begin{itemize}
\item
  \emph{Personality:} tone + style
\item
  \emph{Prompt template:} instructions, constraints
\item
  \emph{Data sources:} files, knowledge bases
\item
  \emph{Actions:} external tools
\end{itemize}

\textbf{Use cases:}

\begin{itemize}
\item
  Automating bioinformatics reports
\item
  Turning handwritten notes into Google Docs
\item
  Cleaning up metadata
\end{itemize}

\section{Personality}\label{personality}

Customize how your chatbot interacts:

\begin{itemize}
\tightlist
\item
  Confidence: a) assured, b) measured
\item
  Amicability: a)friendly, b) neutral
\item
  Professionalism: a) formal, b) casual
\item
  Interactivity: a) engaging, b) informative
\item
  Transparency: a) adaptive, b) consistent
\item
  Lexicography: a) specialized, b) universal
\end{itemize}

\section{Structuring Input \& Output}\label{structuring-input-output}

\subsection{Input Pattern}\label{input-pattern}

Prompt instructs GPT to ask for missing info.

\textbf{Example:}

``Before writing, ask me for the dataset description.''

\subsection{Structured Output (tables, checklists,
format)**}\label{structured-output-tables-checklists-format}

Makes results more usable.

\textbf{Example (table):}

\textbar{} Step \textbar{} Description \textbar{} Tools \textbar{}

\textbar------\textbar-------------\textbar-------\textbar{}

\textbar{} 1 \textbar{} QC \textbar{} FastQC \textbar{}

\textbar{} 2 \textbar{} Alignment \textbar{} BWA \textbar{}

Use markdown formatting options (headers, tables, links, footnotes)!

This is useful when prompting for \textbf{reusable}, \textbf{clean},
\textbf{publication-ready} output.

\section{Tabular data analysis}\label{tabular-data-analysis}

LLMs can understand tables if provided clearly.

\textbf{Example:}

``Using the table below, predict the missing column `Cell type'.''

They can also create synthetic datasets to help test pipelines.

\section{In-Context Learning with Tabular
Data}\label{in-context-learning-with-tabular-data}

Provide columns + examples → model deduces hidden rules.

\textbf{Example:}

``Given columns A and B, predict column C (protein class) based on
observed patterns.''

\section{Synthetic Data Generation}\label{synthetic-data-generation}

Bootstrap ML experiments or testing.

\textbf{Example:}

``Generate a synthetic metadata table of 50 samples with fields:
donor\_age, diagnosis, cell\_type, library\_size.''

\section{ML with Prompts (Your Course
Notes)}\label{ml-with-prompts-your-course-notes}

\textbf{Classification via Prompting}

\begin{itemize}
\tightlist
\item
  Good for simple yes/no or category labeling.
\end{itemize}

\textbf{Example:}

``Label each comment as `informative', `spam', or `uncertain'.''

\textbf{Clustering via Prompting}

\begin{itemize}
\tightlist
\item
  Model groups items into meaningful clusters.
\end{itemize}

\textbf{Example:}

``Cluster the following 30 job descriptions into 5 roles and output an
ASCII taxonomy tree.''

\textbf{Prediction via Prompting}

\begin{itemize}
\tightlist
\item
  Model extrapolates patterns.
\end{itemize}

\textbf{Example:}

``Given the engagement metrics for existing workshop titles, propose 3
new course titles likely to perform well.''

\section{Benchmarking GPT (Your ``Test the GPT''
Notes)}\label{benchmarking-gpt-your-test-the-gpt-notes}

Things to test:

\begin{itemize}
\tightlist
\item
  Factual accuracy
\item
  Reasoning
\item
  Creativity
\item
  Cultural/linguistic alignment
\item
  Output structure
\item
  Handling of adversarial or ambiguous input
\end{itemize}

\textbf{Example:}

``Test whether the model stays consistent when given conflicting
instructions.''

Built your own testerGPT to benchmark other customGPTs!

\subsection{Reliability \& Safety}\label{reliability-safety}

\textbf{Preventing hallucinations}

\begin{itemize}
\tightlist
\item
  Provide all key information first.
\item
  Ask user to confirm before acting.
\end{itemize}

\textbf{Escape Valves} to avoid forced answers.

\textbf{Example:}

``If missing, respond with \textless N/A\textgreater.''

\textbf{Boundaries}

\begin{itemize}
\tightlist
\item
  For safety, compliance, correctness.
\item
  Useful in regulated areas (e.g., healthcare, legal, institutional
  guidelines).
\end{itemize}

\textbf{Example:}

``Do not provide clinical guidance. If asked, offer general information
and suggest contacting a medical professional.''

\textbf{Ambiguity Handling}

\begin{enumerate}
\def\labelenumi{\arabic{enumi}.}
\tightlist
\item
  Identify it
\item
  Ask for clarification or fall back on \textless N/A\textgreater{}
\end{enumerate}

\subsection{Conflict Resolution}\label{conflict-resolution}

When documents/ inputs conflict:

\begin{itemize}
\tightlist
\item
  Use newest
\item
  If one contains the needed topic and the other not, use the relevant
  one
\item
  State uncertainty explicitly
\end{itemize}

\subsection{Adversarial Input
Prevention}\label{adversarial-input-prevention}

\begin{itemize}
\tightlist
\item
  Add railguards
\item
  Disable functionality beyond scope
\item
  Avoid being tricked into harmful or nonsensical replies
\end{itemize}

\bookmarksetup{startatroot}

\chapter{Data analysis with ChatGTP}\label{data-analysis-with-chatgtp}

My notes from the Coursera courses \emph{Advanced Data Analysis with
ChatGPT} and \emph{Building Trustworthy AI} (Code Interpreter) by Jules
White. Digitized from handwritten material via mobile photos and
organized into a coherent structure using ChatGPT (customGPT: ``Notes
Organizer'').

\textbf{Synopsis:} How to use ChatGPT effectively for data analysis,
document processing, workflow planning, and creative problem-solving,
and outline key principles, tools, and frameworks (such as ACHIEVE and
ETAC) that help AI augment human reasoning while maintaining accuracy
and control.

\begin{center}\rule{0.5\linewidth}{0.5pt}\end{center}

\section{General Principles of Working with generative
AI}\label{general-principles-of-working-with-generative-ai}

\textbf{Fundamentals}

\begin{itemize}
\tightlist
\item
  \emph{``LLMs do not generate answers, LLMs generate text. Some of the
  text might provide useful answers''(Jules White)}
\item
  LLM output has an inherent random component (LLMs predict, and the
  same prompt may not provide the exact same output)
\item
  Using generative AI implys taking ownership of the output.
\end{itemize}

\textbf{Treat Code Interpreter as an internal assistant}

Think of the Code Interpreter as a smart intern that needs clear
instructions:

\begin{enumerate}
\def\labelenumi{\arabic{enumi}.}
\tightlist
\item
  \textbf{Break tasks down step-by-step.}
\item
  \textbf{Provide all necessary context and information upfront} -- data
  files, constraints, goals, examples.
\item
  Ask for \textbf{detailed feedback}, e.g.~``What could be improved
  about this figure?''
\item
  \textbf{Don't expect perfection.} Plan to iterate with it.
\item
  \textbf{Always review results} for accuracy, plausibility, and errors.
\item
  Be \textbf{explicit} about:

  \begin{itemize}
  \tightlist
  \item
    What you want to achieve.
  \item
    How you want results presented (tables, plots, summaries, slides).
  \item
    Any constraints (runtime, file formats, colour schemes, etc.).
  \end{itemize}
\item
  \textbf{Provide escape valves}
\end{enumerate}

\begin{center}\rule{0.5\linewidth}{0.5pt}\end{center}

\section{Working with Data Files, Documents, Media and Zip
Files}\label{working-with-data-files-documents-media-and-zip-files}

\subsection{Working with Data Files and
Visualizations}\label{working-with-data-files-and-visualizations}

\subsubsection{Uploading and exploring
data}\label{uploading-and-exploring-data}

Typical workflow with Code Interpreter:

\begin{itemize}
\tightlist
\item
  Upload a dataset (csv) and ask:

  \begin{itemize}
  \tightlist
  \item
    ``Read this file and explain the structure of the data.''
  \item
    ``Show me the columns, data types, and basic statistics.''
  \end{itemize}
\item
  Ask for \textbf{interesting visualizations}:

  \begin{itemize}
  \tightlist
  \item
    ``Create several interesting plots that could be useful for
    exploratory analysis.''
  \item
    ``Try a histogram, a boxplot, and a scatterplot matrix.''
  \end{itemize}
\end{itemize}

You can then iterate:

\begin{itemize}
\tightlist
\item
  ``Refine the labels and axes.''
\item
  ``Use a harmonized colour scheme.''
\item
  ``Order categories in a meaningful way.''
\end{itemize}

\subsubsection{Adding narrative text and
presentations}\label{adding-narrative-text-and-presentations}

You can ask ChatGPT to help with presentation-ready outputs:

\begin{itemize}
\tightlist
\item
  ``Insert each visualization into a slide and write a short caption
  explaining what it shows.''
\item
  ``Update my existing slides with a short narrative for each plot.''
\item
  ``Use consistent terminology and colour coding across all slides.''
\end{itemize}

Example prompt:

\begin{quote}
``Create a one-slide summary for each of these plots, explaining what
the viewer should take away. Use clear language suitable for a
multidisciplinary biomedical audience.''
\end{quote}

\subsection{Working with Documents}\label{working-with-documents}

\subsubsection{Documents as objects vs.~documents as
knowledge}\label{documents-as-objects-vs.-documents-as-knowledge}

Two levels of interaction with documents:

\begin{enumerate}
\def\labelenumi{\arabic{enumi}.}
\tightlist
\item
  \textbf{Document as an object}

  \begin{itemize}
  \tightlist
  \item
    Operations: rename, move, download, edit, zip/unzip.
  \end{itemize}
\item
  \textbf{Document as knowledge}

  \begin{itemize}
  \tightlist
  \item
    Operations: analyze, summarize, outline, extract key statements,
    compare sections, etc.
  \end{itemize}
\end{enumerate}

Being aware of this distinction helps you design useful prompts.

\subsubsection{Useful document prompts}\label{useful-document-prompts}

Examples:

\begin{itemize}
\tightlist
\item
  ``Extract this PDF to plain text and tell me what it contains at a
  high level.''
\item
  ``Reread the text and break it down into the smallest meaningful
  sections with headings.''
\item
  ``Create an outline of the document.''
\item
  ``Identify contradictions, gaps, or redundancies.''
\item
  ``Highlight all places where inclusion criteria are mentioned and
  summarise them.''
\end{itemize}

This improves:

\begin{itemize}
\tightlist
\item
  \textbf{Reasoning} -- you understand the structure more clearly.
\item
  \textbf{Verification} -- easier to check that all key information is
  captured.
\item
  \textbf{Navigation} -- easier to find where specific concepts live.
\end{itemize}

\subsection{Using Canvas in ChatGPT}\label{using-canvas-in-chatgpt}

Canvas is useful for interactive editing and document development.

\subsubsection{Basic Canvas workflow}\label{basic-canvas-workflow}

\begin{enumerate}
\def\labelenumi{\arabic{enumi}.}
\tightlist
\item
  Open \textbf{Canvas}.
\item
  Paste or import your text (or upload a document to extract text).
\item
  Enter a prompt to modify the text or to ask for comments/suggestions.
\item
  Edit the text manually or accept/reject suggested changes.
\item
  Optionally upload additional documents for context (e.g.~guidelines,
  previous drafts).
\item
  Enter more prompts to adjust \textbf{tone}, \textbf{clarity},
  \textbf{length}, or \textbf{structure}.
\item
  Use the Canvas controls (e.g.~icons in the lower corner) to adjust
  text length and polish.
\end{enumerate}

\subsubsection{Two main Canvas entry
paths}\label{two-main-canvas-entry-paths}

\begin{itemize}
\tightlist
\item
  \textbf{Upload a document + prompt}\\
  \textgreater{} ``Extract the text from this document into Canvas and
  clean up formatting.''
\item
  \textbf{Copy \& paste text directly}\\
  Then iterate with prompts to refine.
\end{itemize}

\subsection{Structured Data Workflows}\label{structured-data-workflows}

\subsubsection{Understanding tables}\label{understanding-tables}

Useful prompts:

\begin{itemize}
\tightlist
\item
  ``Read and explain the structure of this table.''
\item
  ``Describe each column, including data type and potential problems.''
\item
  ``Identify missing values, inconsistent labels, and outliers.''
\end{itemize}

Then:

\begin{itemize}
\tightlist
\item
  ``Calculate derived variables step-by-step and show intermediate
  tables.''
\item
  ``Visualise the relationship between variable X and Y, including
  confidence intervals.''
\end{itemize}

\subsubsection{Example workflow}\label{example-workflow}

\begin{enumerate}
\def\labelenumi{\arabic{enumi}.}
\tightlist
\item
  Load dataset (CSV, TSV, Excel, etc.).
\item
  Ask for an overview (dimensions, column types, missingness).
\item
  Clean: standardise categories, handle missing values.
\item
  Derive new variables (ratios, differences, grouping codes).
\item
  Visualise key relationships.
\item
  Summarise results in natural language.
\end{enumerate}

\subsection{Working with Media and Zip
Archives}\label{working-with-media-and-zip-archives}

\subsubsection{Video workflows}\label{video-workflows}

Typical operations:

\begin{itemize}
\tightlist
\item
  ``Extract N frames from this video, evenly spaced.''
\item
  ``Display the extracted frames.''
\item
  ``Resize each image to 300 px wide while preserving aspect ratio.''
\item
  ``Convert to grayscale and increase contrast by 30\%.''
\end{itemize}

This is useful for:

\begin{itemize}
\tightlist
\item
  Creating thumbnails.
\item
  Quick quality checks.
\item
  Preparing training data for image models.
\end{itemize}

\subsubsection{Zip archives}\label{zip-archives}

Zip archives are handy for batch operations:

\begin{itemize}
\tightlist
\item
  Zip multiple outputs into a single file for download.
\item
  Ask ChatGPT to \textbf{propose a folder structure} and \textbf{better
  filenames}.
\end{itemize}

Example prompt:

\begin{quote}
``Resize all images to 300 px width, rename them systematically, and put
them into a zip file for download. Suggest a clear naming scheme based
on sample ID and time point.''
\end{quote}

\section{Workflows and Planning}\label{workflows-and-planning}

\section{Check: Bringing Information into the
Conversation}\label{check-bringing-information-into-the-conversation}

Useful operations:

\begin{itemize}
\tightlist
\item
  ``Read this long text.''
\item
  ``Summarize it in X bullet points.''
\item
  ``Present the key information back to me in a way that I can quickly
  verify.''
\item
  ``Create an index of where specific information (e.g.~inclusion
  criteria, endpoints) can be found.''
\end{itemize}

This helps to:

\begin{itemize}
\tightlist
\item
  Improve reasoning.
\item
  Support verification and auditing.
\item
  Create better documentation and overviews.
\end{itemize}

\subsection{The ETAC Framework for workflows Human + AI Process
Planning}\label{the-etac-framework-for-workflows-human-ai-process-planning}

Think of the overall human--AI workflow as:

\textbf{E -- Extract → T -- Transform → A -- AI / Analyze → C -- Create}

\subsubsection{Extract: modalities and
operations}\label{extract-modalities-and-operations}

\begin{itemize}
\tightlist
\item
  \textbf{Documents (docx, pdf, txt)}

  \begin{itemize}
  \tightlist
  \item
    Text extraction
  \item
    Metadata extraction
  \item
    Entity recognition
  \end{itemize}
\item
  \textbf{Video (mp4, avi)}

  \begin{itemize}
  \tightlist
  \item
    Frame extraction
  \item
    Object detection
  \item
    (Sometimes) face recognition
  \end{itemize}
\item
  \textbf{Audio (mp3, wav)}

  \begin{itemize}
  \tightlist
  \item
    Speech-to-text
  \item
    Noise reduction
  \item
    Feature extraction
  \end{itemize}
\item
  \textbf{Datasets (csv, tsv, json, xlsx)}

  \begin{itemize}
  \tightlist
  \item
    Load data
  \item
    Preprocess
  \item
    Feature engineering
  \end{itemize}
\item
  \textbf{Images (jpg, png, etc.)}

  \begin{itemize}
  \tightlist
  \item
    Object detection
  \item
    Segmentation
  \item
    Feature extraction
  \end{itemize}
\end{itemize}

\subsubsection{Transform: data and
images}\label{transform-data-and-images}

\begin{itemize}
\tightlist
\item
  \textbf{Structured data}

  \begin{itemize}
  \tightlist
  \item
    Clean and normalise
  \item
    Drop or merge columns
  \item
    Recode values
  \item
    Compute derived features
  \end{itemize}
\item
  \textbf{Images}

  \begin{itemize}
  \tightlist
  \item
    Subselect regions (ROI)
  \item
    Crop and resize
  \item
    Adjust contrast, brightness, etc.
  \end{itemize}
\end{itemize}

Always prefer \textbf{small, controlled steps} that you can undo or
repeat.

\subsection{Analyze: From raw material to
insight}\label{analyze-from-raw-material-to-insight}

Analysis-related operations include:

\begin{itemize}
\tightlist
\item
  Find ambiguities and inconsistencies.
\item
  Analyse and summarise content.
\item
  Raise clarifying questions for the user.
\item
  Generate examples and counter-examples.
\end{itemize}

You can also maintain a ``library'' or mental catalogue of:

\begin{itemize}
\tightlist
\item
  File types
\item
  Common operations
\item
  Good default workflows
\end{itemize}

\subsection{Planning Workflows -- Human and AI
Roles}\label{planning-workflows-human-and-ai-roles}

\subsubsection{Human planning}\label{human-planning}

\begin{itemize}
\tightlist
\item
  Lay out the process manually first:

  \begin{itemize}
  \tightlist
  \item
    Inputs
  \item
    Intermediate artefacts
  \item
    Outputs
  \end{itemize}
\item
  Decide what you want to control vs.~what AI can suggest.
\end{itemize}

\subsubsection{AI planning}\label{ai-planning}

Ask ChatGPT:

\begin{quote}
``Propose a detailed step-by-step plan for this task.''
\end{quote}

Then refine:

\begin{itemize}
\item
  Expand each bullet into \textbf{2--3 actionable substeps}.
\item
  Ensure each step is something:

  \begin{itemize}
  \tightlist
  \item
    You can perform, or
  \item
    The Python environment can perform.
  \end{itemize}
\item
  Turn the result into an \textbf{actionable plan} or checklist.
\item
  Verify that the sequence of step is executable
\end{itemize}

\subsection{AI Planning Techniques and
Iteration}\label{ai-planning-techniques-and-iteration}

Good practice when executing a plan with Code Interpreter:

\begin{enumerate}
\def\labelenumi{\arabic{enumi}.}
\tightlist
\item
  Execute steps \textbf{one by one}.
\item
  Save results of each step into \textbf{separate files} (for
  reproducibility and debugging).
\item
  At the beginning of each step:

  \begin{itemize}
  \tightlist
  \item
    Reflect on prior output.
  \item
    Check alignment with goals.
  \end{itemize}
\item
  If something breaks:

  \begin{itemize}
  \tightlist
  \item
    Stop, diagnose, fix the step.
  \item
    Then continue.
  \end{itemize}
\end{enumerate}

You can also:

\begin{itemize}
\tightlist
\item
  Store the plan itself in a separate file for later reuse.
\item
  Use ``flipped interaction'': ask AI to improve \emph{your} plan,
  rather than just proposing its own.
\end{itemize}

\subsection{Alternative Approaches and Tool
Choice}\label{alternative-approaches-and-tool-choice}

Not all plans will work on the first try:

\begin{itemize}
\tightlist
\item
  Some approaches will fail due to limitations of the model or tools.
\item
  It's valuable to \textbf{try alternative methods} and compare.
\end{itemize}

Tool choice matters:

\begin{itemize}
\tightlist
\item
  \textbf{GPT / ChatGPT} -- reasoning, summarising, planning, generating
  text.
\item
  \textbf{Python / other tools} -- computation, heavy data processing,
  integration with external systems.
\end{itemize}

Often you manually choose \textbf{when} to ask the LLM to ``do
something'' and when to switch to specialised tools.

\subsection{Fixing and Spotting
Errors}\label{fixing-and-spotting-errors}

\subsubsection{Fixing errors}\label{fixing-errors}

\begin{itemize}
\tightlist
\item
  Go back to the prompt \textbf{before} the error occurred and edit it.
\item
  Number steps in your plan so you can say: \textgreater{} ``Repeat step
  4, but show the intermediate table before filtering.''
\item
  Ask for explicit error reporting and logging.
\end{itemize}

\subsubsection{Spotting hallucinations}\label{spotting-hallucinations}

Warning signs:

\begin{itemize}
\tightlist
\item
  \textbf{Input/output ratio is very small} -- tiny context, huge
  confident answer.
\item
  A lot of \textbf{artificial data is generated} without clear reason.
\item
  Long conversation history with many state changes.
\end{itemize}

Additional caution when:

\begin{itemize}
\tightlist
\item
  Using conversation history as implicit input.
\item
  Re-running code cells on modified data without re-checking
  assumptions.
\end{itemize}

\subsection{Consistency, Test Cases, and Source
Linking}\label{consistency-test-cases-and-source-linking}

To build more trustworthy workflows:

\begin{itemize}
\tightlist
\item
  Use identifiers and citations:

  \begin{itemize}
  \tightlist
  \item
    Page numbers, line numbers, section IDs.
  \end{itemize}
\item
  Preserve references to \textbf{original sources} when summarising.
\item
  Build \textbf{test cases}:

  \begin{itemize}
  \tightlist
  \item
    Use synthetic data to test specific actions/conditions.
  \item
    Check whether the workflow behaves as expected.
  \end{itemize}
\end{itemize}

\subsection{Breaking Down and Assembling Larger
Projects}\label{breaking-down-and-assembling-larger-projects}

\subsubsection{Breaking down a large
project}\label{breaking-down-a-large-project}

\begin{itemize}
\tightlist
\item
  Start by \textbf{creating an index}:

  \begin{itemize}
  \tightlist
  \item
    Sections, topics, questions.
  \end{itemize}
\item
  Then use the index to drive operations:

  \begin{itemize}
  \tightlist
  \item
    ``Using this index, extract all content related to X.''
  \end{itemize}
\end{itemize}

\subsubsection{Assembling output into a larger
product}\label{assembling-output-into-a-larger-product}

Incremental workflow:

\begin{enumerate}
\def\labelenumi{\arabic{enumi}.}
\tightlist
\item
  Run analyses chunk by chunk.
\item
  Produce \textbf{individual output files} (markdown, figures, tables).
\item
  At the end, combine them into:

  \begin{itemize}
  \tightlist
  \item
    A report
  \item
    A slide deck
  \item
    A book chapter
  \end{itemize}
\end{enumerate}

This is particularly useful for \textbf{text analysis}:

\begin{itemize}
\tightlist
\item
  First build an outline.
\item
  Then expand each section via incremental analysis.
\item
  Finally assemble using the outline as a backbone.
\end{itemize}

\subsection{Turning Conversations into
Software}\label{turning-conversations-into-software}

You can move from interactive exploration to reusable tools. (From
workflow to program)

Prompt example:

\begin{quote}
``Turn this process into a Python program that I can download and run
locally. Include clear comments and command-line arguments.''
\end{quote}

Good practice:

\begin{itemize}
\tightlist
\item
  Inspect the code \textbf{carefully} before running it.
\item
  Prefer running in a \textbf{safe environment} (sandbox, container,
  virtualenv).
\item
  Keep configuration (file paths, parameters) outside the core logic
  where possible.
\end{itemize}

\section{Problem Selection and Trustworthy
AI}\label{problem-selection-and-trustworthy-ai}

\subsection{Picking Good Problems for AI (ACHIEVE-aligned
view)}\label{picking-good-problems-for-ai-achieve-aligned-view}

Good AI candidates:

\begin{itemize}
\tightlist
\item
  Easy to \textbf{check} whether the answer is correct.
\item
  A \textbf{partial solution} already saves substantial time.
\item
  The task benefits from \textbf{alternative perspectives} or
  reformulations.
\item
  The workflow can be decomposed into clear steps with observable
  outputs.
\end{itemize}

AI should \textbf{augment}, not replace, human judgment, creativity or
intelligence.

\subsection{The ACHIEVE Framework
(Overview)}\label{the-achieve-framework-overview}

ACHIEVE is a way to think about how AI can augment human intelligence:

\begin{itemize}
\tightlist
\item
  \textbf{A -- Aid human coordination}
\item
  \textbf{C -- Cut tedium / Reduce tedious work}
\item
  \textbf{H -- Provide a safety net (spot errors/risks)}
\item
  \textbf{I -- Inspire creativity}
\item
  \textbf{E -- Execute and scale ideas}
\item
  \textbf{V -- Visibility and verification}
\item
  \textbf{E -- Experimentation and exploration}
\end{itemize}

Below, concrete examples and use cases.

\subsection{ACHIEVE -- Detailed Use
Cases}\label{achieve-detailed-use-cases}

\subsubsection{Aid human coordination}\label{aid-human-coordination}

Example: planning a workshop series.

\begin{itemize}
\tightlist
\item
  Summarise meeting notes into key decisions and actions.
\item
  Identify actors and assign responsibilities.
\item
  Spot ambiguities and potential conflicts.
\item
  Propose clarifications and conflict-resolution options.
\item
  Draft follow-up emails and reminders.
\end{itemize}

\subsubsection{Cut tedium / Reduce tedious
work}\label{cut-tedium-reduce-tedious-work}

Example: cleaning data from a registration form.

\begin{itemize}
\tightlist
\item
  Map similar free-text entries onto consistent categories.
\item
  Create data dictionaries and lookup tables.
\item
  Group related terms and suggest standardised labels.
\item
  Generate code snippets for repetitive cleaning operations.
\item
  Visualise distributions to check for odd patterns.
\end{itemize}

\subsubsection{Help provide a safety
net}\label{help-provide-a-safety-net}

Example: travel reimbursement workflows.

\begin{itemize}
\tightlist
\item
  Check forms for completeness and consistency.
\item
  Compare entries against internal rules or policies.
\item
  Highlight potential issues or missing information.
\item
  Suggest corrections or follow-up questions.
\item
  Important: \textbf{final approval remains human.}
\end{itemize}

\subsubsection{Inspire creativity}\label{inspire-creativity}

Example: designing a workshop.

\begin{itemize}
\tightlist
\item
  Ask it questions about your content to rethink structure:

  \begin{itemize}
  \tightlist
  \item
    ``What natural modules do you see in this outline?''
  \end{itemize}
\item
  Generate follow-up or bonus content ideas.
\item
  Suggest logistic setups and communication strategies.
\item
  Design templates for slides, handouts, and emails.
\item
  Help develop ideas into more detailed formats (e.g.~exercises, case
  studies).
\end{itemize}

\subsubsection{Execute and scale}\label{execute-and-scale}

\begin{itemize}
\tightlist
\item
  Turn proven workflows into reusable templates.
\item
  Create ``starter kits'' (data + notebook + documentation).
\item
  Generate parameterised scripts for common tasks.
\item
  Build small internal tools that others can run with minimal effort.
\end{itemize}

\subsubsection{Visibility and
verification}\label{visibility-and-verification}

\begin{itemize}
\tightlist
\item
  Create diagrams of workflows and data flows.
\item
  Summarise long documents and highlight key decisions.
\item
  Create indexes and cross-reference tables.
\item
  Keep a clear link between outputs and their sources (citations, IDs).
\end{itemize}

\subsubsection{Experimentation and
exploration}\label{experimentation-and-exploration}

\begin{itemize}
\tightlist
\item
  Quickly test multiple approaches (alternative prompts, models,
  parameters).
\item
  Use synthetic data to explore edge cases safely.
\item
  Compare strategies side-by-side and reflect on trade-offs.
\end{itemize}

\subsection{ACHIEVE for Augmented Intelligence
(Recap)}\label{achieve-for-augmented-intelligence-recap}

The goal is to \textbf{augment and amplify} human creativity and
problem-solving, not replace people.

ACHIEVE helps you:

\begin{itemize}
\tightlist
\item
  Spend more time on high-value, human tasks.
\item
  Use AI as a thinking partner and productivity booster.
\item
  Support communication, alignment, and shared understanding.
\end{itemize}

\subsubsection{Aid human communications}\label{aid-human-communications}

\begin{itemize}
\tightlist
\item
  Draft email templates.
\item
  Summarise meetings and highlight open questions.
\item
  Clarify ambiguities and identify conflicts.
\item
  Map actors and actions across a project.
\end{itemize}

\subsubsection{Cut tedious tasks}\label{cut-tedious-tasks}

\begin{itemize}
\tightlist
\item
  Automate repetitive formatting and cleaning.
\item
  Batch-generate routine outputs.
\item
  Scaffold code and documentation.
\end{itemize}

\subsubsection{Help privide safety net}\label{help-privide-safety-net}

AI can:

\begin{itemize}
\tightlist
\item
  Highlight ambiguities and potential conflict points.
\item
  Surface missing information or inconsistent logic.
\item
  Help evaluate risks and suggest mitigations.
\end{itemize}

But:

\begin{itemize}
\tightlist
\item
  Humans remain responsible for final decisions.
\end{itemize}

\subsubsection{Ideation}\label{ideation}

\begin{itemize}
\tightlist
\item
  Generate many ideas quickly and cheaply.
\item
  Use AI for \textbf{low-risk brainstorming} and exploration.
\item
  Helpful for diagramming and visualising possibilities.
\item
  The human user:

  \begin{itemize}
  \tightlist
  \item
    Selects promising ideas.
  \item
    Combines and adapts them.
  \item
    Maintains ownership and direction of the process.
  \end{itemize}
\end{itemize}

\subsubsection{Extract and Filtering}\label{extract-and-filtering}

\begin{itemize}
\tightlist
\item
  Filter context information so reasoning stays tractable.
\item
  Summarise or analyse while keeping references to original sources.
\item
  Ensure filtering is \textbf{immediately checkable} (quote, cite,
  link).
\end{itemize}

\subsection{Building Trustworthy AI}\label{building-trustworthy-ai}

Dimensions of trustworthiness:

\begin{itemize}
\tightlist
\item
  \textbf{Veracity} -- conformity to facts.
\item
  \textbf{Accuracy} -- numerical / logical correctness.
\item
  \textbf{Cost of checking} vs \textbf{cost of producing} the answer.
\end{itemize}

Remember:

\begin{itemize}
\tightlist
\item
  Generative AI is \textbf{not} a primary source of facts.
\item
  Treat outputs as \textbf{drafts to be checked}, not final truth.
\end{itemize}

Good practices:

\begin{itemize}
\tightlist
\item
  Mark outputs by certainty or evidence level.
\item
  Use visual aids (flow diagrams, decision trees) to help reasoning.
\item
  Apply risk assessments, especially for:

  \begin{itemize}
  \tightlist
  \item
    Physical risks
  \item
    Social / psychological risks
  \item
    Economic or career impacts
  \item
    Data protection and privacy
  \end{itemize}
\end{itemize}

AI can: - Point you to relevant sections in documents (tracing) -
Maintain a ``trail of generation'' -- where information came from. -
Spot missing or fragmented documentation and help mitigate consequences.

This is \textbf{especially useful with sensitive data}, where you must
know exactly:

\begin{itemize}
\tightlist
\item
  What was used.
\item
  Why it was used.
\item
  How outputs were derived.
\end{itemize}

Remember the importance of expertise:

\begin{itemize}
\tightlist
\item
  Appropriate \textbf{domain expertise} is crucial to interpret outputs.
\item
  AI cannot replace expertise; it can only \textbf{amplify} it.
\item
  In biomedical and clinical research, expert review is non-negotiable.
\end{itemize}




\end{document}
